Joint optimization of signal control and road pricing has a bilevel programming tradition dating
to at least \citet{smith1997}, who set signal timings and road prices together within a link flow
framework. \citet{clegg2001} extend this to a multi modal elastic demand model, tested on a small
network. \citet{mao2013} model signal settings and congestion tolls as a bilevel program.
\citet{smith2022} show that using link exit prices in a combined control policy maximizes network
capacity with zero queues in steady state. No source located in this review reports a measured
combined effect tested against the sum of each lever applied alone, the question H3 asks.

Prior Kathmandu specific traffic modeling is limited to single intersections. \citet{tamang2026}
model the Lainchaur intersection in VISSIM, finding major approaches at Level of Service F.
\citet{tiwari2024satdobato} optimize signals at Satdobato and Gwarko in SIDRA. \citet{luitel2025}
review the city's congestion without modeling an intervention, and \citet{bhatta2024} estimates a
congestion cost for one road section. No located study models the city at network scale or tests
pricing or demand management there.
